\begingroup
\renewcommand{\arraystretch}{1.35}
\InterpolationArticle{3}
Es sei $X$ eine Funktion der Unbestimmten $x$ von der Form
\[
\alpha+\beta x+\gamma xx+\delta x^3+\text{usw.}
\]
die nicht bis ins Unendliche fortgesetzt wird, sondern abbricht und nicht über das Glied mit $x^{m-1}$ hinausgeht, sodass die Anzahl der Koeffizienten nicht größer als $m$ ist. Sind dann für $m$ verschiedene Werte von $x$, nämlich $a,b,c,d\ldots$, die entsprechenden Werte von $X$ bekannt, nämlich $=A,B,C,D\ldots$, so lässt sich daraus der Wert von $X$, der einem beliebigen anderen Wert von $x$ entspricht, auf folgende Weise übersichtlich ermitteln. Es sei $t$ der neue Wert von $x$ und $T$ der entsprechende Wert von $X$, sodass die folgenden $m+1$ Gleichungen gelten
\[
\begin{aligned}
A&=\alpha+\beta a+\gamma aa+\delta a^3+\ldots\\
B&=\alpha+\beta b+\gamma bb+\delta b^3+\ldots\\
C&=\alpha+\beta c+\gamma cc+\delta c^3+\ldots\\
D&=\alpha+\beta d+\gamma dd+\delta d^3+\ldots\\
&\text{usw.}\\
T&=\alpha+\beta t+\gamma tt+\delta t^3+\ldots
\end{aligned}
\]
\endgroup
\par
\begingroup
\renewcommand{\arraystretch}{1.45}
\SourceLeaf{282}
Man multipliziere diese Gleichungen der Reihe nach mit
\[
\begin{gathered}
\frac{1}{(a-b)(a-c)(a-d)\mathinner{.\,.\,.\,.}(a-t)}\\[4pt]
\frac{1}{(b-a)(b-c)(b-d)\mathinner{.\,.\,.\,.}(b-t)}\\[4pt]
\frac{1}{(c-a)(c-b)(c-d)\mathinner{.\,.\,.\,.}(c-t)}\\[4pt]
\frac{1}{(d-a)(d-b)(d-c)\mathinner{.\,.\,.\,.}(d-t)}\\
\text{usw.}\\
\frac{1}{(t-a)(t-b)(t-c)\mathinner{.\,.\,.\,.}}
\end{gathered}
\]
und addiere die Produkte, sodass sich ergibt
\[
\begin{aligned}
&\frac{A}{(a-b)(a-c)(a-d)\mathinner{.\,.\,.}(a-t)}\\
&{}+\frac{B}{(b-a)(b-c)(b-d)\mathinner{.\,.\,.}(b-t)}\\
&{}+\frac{C}{(c-a)(c-b)(c-d)\mathinner{.\,.\,.}(c-t)}\\
&{}+\frac{D}{(d-a)(d-b)(d-c)\mathinner{.\,.\,.}(d-t)}\\
&{}+\text{usw.}\\
&{}+\frac{T}{(t-a)(t-b)(t-c)(t-d)\mathinner{.\,.\,.}}=W
\end{aligned}
\]
Aus Artikel 1, wo $m$ dasselbe bezeichnete, was für uns hier $m+1$ ist, folgt dann leicht $W=0$; durch Multiplikation mit $(t-a)(t-b)(t-c)\mathinner{.\,.\,.}$ ergibt sich daher
\[
\begin{aligned}
T={}&\frac{(t-b)(t-c)(t-d)\mathinner{.\,.\,.}}{(a-b)(a-c)(a-d)\mathinner{.\,.\,.}}A\\
&{}+\frac{(t-a)(t-c)(t-d)\mathinner{.\,.\,.}}{(b-a)(b-c)(b-d)\mathinner{.\,.\,.}}B\\
&{}+\frac{(t-a)(t-b)(t-d)\mathinner{.\,.\,.}}{(c-a)(c-b)(c-d)\mathinner{.\,.\,.}}C\\
&{}+\frac{(t-a)(t-b)(t-c)\mathinner{.\,.\,.}}{(d-a)(d-b)(d-c)\mathinner{.\,.\,.}}D\\
&{}+\text{usw.}
\end{aligned}
\]
\InterpolationArticle{4}
Die im vorhergehenden Artikel gefundene Formel ist so beschaffen, dass ohne jede Rechnung unmittelbar erkennbar ist: Wird für $t$ eine der Größen $a,b,c,d\mathinner{.\,.}$ eingesetzt, so ergibt sich der
entsprechende Wert $A,B,C,D\mathinner{.\,.}$. Nicht nur in dieser Hinsicht empfiehlt sie sich: Für den praktischen Gebrauch ist sie sicherlich bei Weitem die bequemste, wenigstens wenn nur ein ein\SourceLeaf{283}%
ziger Wert von $X$ aus den gegebenen Werten zu berechnen ist. Sind mehrere
zu ermitteln, so kann mitunter die folgende Umformung unserer Formel
vorzuziehen sein
\[
\begin{aligned}
T={}&A+(t-a)\left(\frac{A}{a-b}+\frac{B}{b-a}\right)\\
&{}+(t-a)(t-b)\left(\frac{A}{(a-b)(a-c)}+\frac{B}{(b-a)(b-c)}+\frac{C}{(c-a)(c-b)}\right)\\
&{}+(t-a)(t-b)(t-c)\left(\frac{A}{(a-b)(a-c)(a-d)}+\frac{B}{(b-a)(b-c)(b-d)}\right.\\
&\qquad\left.{}+\frac{C}{(c-a)(c-b)(c-d)}+\frac{D}{(d-a)(d-b)(d-c)}\right)\\
&{}+\text{usw.}
\end{aligned}
\]
(Sie lässt sich aus der Formel von Artikel 3 sehr leicht herleiten, wenn man sich die Anzahl der Größen
$a,b,c,d\mathinner{.\,.\,.}$ von der einen $a$ zu den beiden $a,b$, dann zu den drei $a,b,c$ usw. allmählich wachsend denkt
und jeden auf diese Weise entstehenden Wert von $T$ vom folgenden subtrahiert.)
Man bezeichne die Koeffizienten wie folgt:
\[
\begin{gathered}\frac{A}{a-b}+\frac{B}{(b-a)},\quad
\frac{A}{(a-b)(a-c)}+\frac{B}{(b-a)(b-c)}+\frac{C}{(c-a)(c-b)}\text{ usw.}\\
=A',\ A''\text{ usw.}\end{gathered}
\]
\Needspace{23\baselineskip}
Ferner bezeichne $B',B''$ usw. das, was aus $A',A''$ usw. entsteht, wenn für $a,b,c\mathinner{.\,.\,.}$
der Reihe nach $b,c,d\mathinner{.\,.\,.}$ und für $A,B,C\mathinner{.\,.\,.}$ der Reihe nach $B,C,D\mathinner{.\,.\,.}$ geschrieben wird; ebenso bezeichne $C',C''$ usw. das, was aus $A',A''$ usw. entsteht, wenn für $a,b,c\mathinner{.\,.\,.}$ der Reihe nach $c,d,e\mathinner{.\,.\,.}$ und für $A,B,C\mathinner{.\,.\,.}$ der Reihe nach $C,D,E\mathinner{.\,.\,.}$ eingesetzt wird, und so weiter. Dann erhalten wir
\[
\begin{array}{lll}
A'=\dfrac{A-B}{a-b},&B'=\dfrac{B-C}{b-c},&C'=\dfrac{C-D}{c-d}\ \mathinner{.\,.\,.}\\[5pt]
A''=\dfrac{A'-B'}{a-c},&B''=\dfrac{B'-C'}{b-d},&C''=\dfrac{C'-D'}{c-e}\ \mathinner{.\,.\,.}\\[5pt]
A'''=\dfrac{A''-B''}{a-d},&B'''=\dfrac{B''-C''}{b-e},&C'''=\dfrac{C''-D''}{c-f}\ \mathinner{.\,.\,.}\\
\text{usw.}&&
\end{array}
\]
und
\[\begin{aligned}
T&=A+A'(t-a)+A''(t-a)(t-b)\\&+A'''(t-a)(t-b)(t-c)+\mathinner{.\,.}
\end{aligned}\]
Die Anzahl der hier zu berechnenden Größen $A',B',C'$ usw. beträgt $m-1$, die Anzahl der
Größen $A'',B'',C''$ usw. beträgt $m-2$ und so weiter.
\InterpolationArticle{5}
Auch die folgende Änderung der Formel von Artikel 3 lässt sich
mit Nutzen anwenden. Es sei $\mathfrak{X}$ eine beliebige ganzrationale Funktion der Unbestimmten $x$, jedoch nicht über den
\SourceLeaf{284}
$m-1^{\mathrm{ten}}$ Grad hinausgehend; ferner seien $\mathfrak{A},\mathfrak{B},\mathfrak{C}.\mathfrak{D}\mathinner{.\,.\,.\,.}\mathfrak{T}$ ihre Werte für $x=a,b,c,d\mathinner{.\,.\,.}t$
der Reihe nach. Dann gilt nach Artikel 1
\[
\begin{aligned}
0={}&\frac{\mathfrak{A}}{(a-b)(a-c)(a-d)\mathinner{.\,.\,.}(a-t)}\\
&{}+\frac{\mathfrak{B}}{(b-a)(b-c)(b-d)\mathinner{.\,.\,.}(b-t)}\\
&{}+\frac{\mathfrak{C}}{(c-a)(c-b)(c-d)\mathinner{.\,.\,.}(c-t)}\\
&{}+\frac{\mathfrak{D}}{(d-a)(d-b)(d-c)\mathinner{.\,.\,.}(d-t)}\\
&{}+\text{usw.}\\
&{}+\frac{\mathfrak{T}}{(t-a)(t-b)(t-c)(t-d)\mathinner{.\,.\,.}}
\end{aligned}
\]
Hieraus lässt sich leicht herleiten
\[
\begin{aligned}
T={}&\mathfrak{T}\\
&{}+\frac{(t-b)(t-c)(t-d)\mathinner{.\,.\,.}}{(a-b)(a-c)(a-d)\mathinner{.\,.\,.}}(A-\mathfrak{A})\\
&{}+\frac{(t-a)(t-c)(t-d)\mathinner{.\,.\,.}}{(b-a)(b-c)(b-d)\mathinner{.\,.\,.}}(B-\mathfrak{B})\\
&{}+\frac{(t-a)(t-b)(t-d)\mathinner{.\,.\,.}}{(c-a)(c-b)(c-d)\mathinner{.\,.\,.}}(C-\mathfrak{C})\\
&{}+\frac{(t-a)(t-b)(t-c)\mathinner{.\,.\,.}}{(d-a)(d-b)(d-c)\mathinner{.\,.\,.}}(D-\mathfrak{D})\\
&{}+\text{usw.}
\end{aligned}
\]
\Needspace{12\baselineskip}
Wird die Funktion $\mathfrak{X}$ so gewählt, dass sie die Werte der Funktion $X$, nämlich $A,B,C,D\mathinner{.\,.\,.}$,
annähernd darstellt, so gewinnen wir den Vorteil, dass die Koeffizienten
$\frac{(t-b)(t-c)(t-d)\mathinner{.\,.\,.}}{(a-b)(a-c)(a-d)\mathinner{.\,.\,.}}$ usw.
nur mit kleinen Größen zu multiplizieren sind. Für $\mathfrak{X}$ kann auch eine konstante Größe
angenommen werden, zum Beispiel eine der Größen $A,B,C,D\mathinner{.\,.\,.}$; in diesem Fall fällt ein Glied aus dem Wert von
$T$ weg. Oder $\mathfrak{X}$ kann so bestimmt werden, dass zwei oder mehr dieser Größen
durch $\mathfrak{X}$ dargestellt werden; dann fallen ebenso viele Glieder aus $T$ weg.
\InterpolationArticle{6}
Das Vorhergehende beruht auf der Voraussetzung, dass die Funktion $X$ nicht über die Potenz
$x^{m-1}$ hinausgeht oder dass die $m^{\mathrm{te}}$ Differenz zusammen mit den höheren Differenzen verschwindet; in
diesem Fall ist die Interpolationsmethode streng richtig. Trifft aber jene Voraussetzung
nicht zu, so zielt die Interpolation nur darauf, anstelle der Funktion $X$ die \textit{einfachste} Funktion zu finden, die den vorgegebenen Werten $A,B,C,D\mathinner{.\,.\,.}$ genügt. Damit wir nun
den Fehler beurteilen können, der durch die vernachlässigten höheren Differenzen entsteht,
\SourceLeaf{285}
seien die Glieder in \(X\) nach der Potenz \(x^{m-1}\) die folgenden: \(\mu x^m+\nu x^{m+1}+\text{usw.}\). Somit gilt in Artikel 3
nicht \(W=0\), sondern, wie aus Artikel 1 folgt, \(W=\mu+\nu(a+b+c+d+\mathinner{.\,.\,.}+t)+\text{usw.}\)
Daher muss zu dem dort angegebenen Wert von \(T\) noch Folgendes hinzugefügt werden
\[
(t-a)(t-b)(t-c)(t-d)\mathinner{.\,.\,.}\times\{\mu+\nu(a+b+c+d+\mathinner{.\,.\,.}+t)+\text{usw.}\}
\]

\endgroup
